\section{Kaczmarz：逐行投影大型线性系统}
\label{sec:08-Numerical-Analysis/05-Numerical-Linear-Algebra/07}

一台 CT 设备每旋转一个角度就发出一束射线，每束射线给出一条方程：沿这条路径的衰减总和等于探测器读数。一次扫描下来有几千万条这样的方程，未知量是几百万个体素。矩阵从来没有被完整写下来过——行是一条一条到达的，而且到达之后就该被丢掉。

上一节的做法全都假设矩阵至少能被乘一次，也就是至少整个在内存里。这里连这一条也不成立。剩下的问题是：一次只看得见一行，还能不能往解的方向走。

\subsection{一次只用一行，就是往一张超平面上投影}

把第 \(i\) 条方程写成 \(a_i^{\trans}x=b_i\)。它在参数空间里画出一张超平面：法向是 \(a_i\)，位置由 \(b_i\) 定。整个系统 \(Ax=b\) 的解，就是所有这些超平面的公共交点。

手上只有第 \(i\) 行时，能做的最保守的一件事是：在满足这一条方程的前提下，尽量少动当前点。写成

\[
x_{k+1}=\argmin_{x}\ \norm{x-x_k}_2^2
\quad\text{s.t.}\quad a_i^{\trans}x=b_i .
\]

Lagrange 条件给出 \(x-x_k=\nu a_i\)，代回约束解出 \(\nu\)，得到

\[
x_{k+1}=x_k+\frac{b_i-a_i^{\trans}x_k}{\norm{a_i}_2^2}\,a_i .
\]

分子是当前点在这条方程上的残差，分母把行向量的长度归掉。这一步之后第 \(i\) 条方程精确成立，下一行的投影又会把它轻微破坏，但只要系统一致，反复投影就会朝公共交点收缩。整个更新只用到 \(a_i\) 和 \(b_i\)：一行数据、一次内积、一次向量加法。矩阵不需要进内存，甚至不需要存在。

\begin{center}
\begin{tikzpicture}[x=1.05cm,y=1.05cm,font=\small,>=stealth]
  \begin{scope}
    \draw[black!55] (0,1.1) -- (4,2.1);
    \draw[black!55] (2.48,0) -- (1.58,3.0);
    \node[right,font=\scriptsize,black!65] at (4.0,2.1) {\(a_1^{\trans}x=b_1\)};
    \node[above,font=\scriptsize,black!65] at (1.58,3.0) {\(a_2^{\trans}x=b_2\)};
    \draw[->,thick,blue!62] (0.40,2.70) -- (0.75,1.29);
    \draw[->,thick,blue!62] (0.75,1.29) -- (1.98,1.66);
    \fill[black!80] (0.40,2.70) circle (1.8pt);
    \node[left,font=\scriptsize,black!75] at (0.36,2.74) {\(x_0\)};
    \fill[blue!62] (0.75,1.29) circle (1.6pt);
    \node[below,font=\scriptsize,blue!70] at (0.75,1.20) {\(x_1\)};
    \node[below,font=\scriptsize,blue!70] at (1.80,1.52) {\(x_2\)};
    \fill[red!70!black] (2.00,1.60) circle (2.0pt);
    \node[right,font=\scriptsize,red!70!black] at (2.16,1.78) {\(x_\star\)};
    \node[font=\footnotesize] at (2.0,-0.60) {两行接近正交};
  \end{scope}
  \begin{scope}[shift={(5.3,0)}]
    \draw[black!55] (0,1.1) -- (4,2.1);
    \draw[black!55] (0,0.5) -- (4,2.7);
    \node[right,font=\scriptsize,black!65] at (4.0,2.1) {\(a_1^{\trans}x=b_1\)};
    \node[right,font=\scriptsize,black!65] at (4.0,2.7) {\(a_2^{\trans}x=b_2\)};
    \draw[->,thick,orange!85!black] (0.50,3.00) -- (0.92,1.33);
    \draw[->,thick,orange!85!black] (0.92,1.33) -- (1.05,1.08);
    \draw[->,thick,orange!85!black] (1.05,1.08) -- (0.99,1.35);
    \draw[->,thick,orange!85!black] (0.99,1.35) -- (1.12,1.11);
    \draw[->,thick,orange!85!black] (1.12,1.11) -- (1.05,1.36);
    \draw[->,thick,orange!85!black] (1.05,1.36) -- (1.17,1.15);
    \fill[black!80] (0.50,3.00) circle (1.8pt);
    \node[left,font=\scriptsize,black!75] at (0.46,3.04) {\(x_0\)};
    \fill[red!70!black] (2.00,1.60) circle (2.0pt);
    \node[right,font=\scriptsize,red!70!black] at (2.10,1.68) {\(x_\star\)};
    \node[font=\footnotesize] at (2.0,-0.60) {两行夹角小};
  \end{scope}
\end{tikzpicture}

\emph{左边两步就到了交点；右边六步还堵在角落里来回蹭——两张超平面的夹角决定每次投影能吃掉多少误差，而这个夹角正是条件数在二维下的样子。}
\end{center}

图里右边那种缓慢的来回，是所有病态系统在 Kaczmarz 下的标准症状。夹角越小，两次连续投影抵消掉的分量越多，净前进越少；夹角趋于零时超平面几乎重合，迭代沿着它们的公共方向几乎不动。\hyperref[sec:08-Numerical-Analysis/01-Floating-Point-and-Error/03]{``条件数衡量问题本身的敏感性''}里那个量在这里有了直接的几何形象：条件数大，就是存在一组行几乎线性相关，也就是存在一个夹角接近零的角落。

\subsection{随机抽行把收敛率变成了一条定理}

固定按 \(1,2,\ldots,m\) 的顺序循环有个隐患：行的排列顺序可能恰好把几张近乎平行的超平面排在一起，让迭代长期困在同一个角落。这类失败与矩阵本身无关，只与你怎么读数据有关，因此也可以只靠改读法解决。

随机 Kaczmarz 按行范数的平方抽样，

\[
P(I_k=i)=\frac{\norm{a_i}_2^2}{\norm{A}_F^2},
\]

于是有期望意义下的线性收敛：

\[
\E\norm{x_k-x_\star}_2^2
\le\Bigl(1-\frac{\sigma_{\min}(A)^2}{\norm{A}_F^2}\Bigr)^{k}\norm{x_0-x_\star}_2^2 .
\]

条件对账要看清楚四件事。这个界要求系统一致且 \(x_\star\) 唯一，噪声数据不适用；它是对随机抽样取期望，单条轨迹可以在某几步上变糟；收敛因子里的 \(\norm{A}_F^2/\sigma_{\min}^2\) 是一个尺度化的条件数平方，病态时它会大到让括号里的收缩因子贴近 1；最后，界与 \(m\) 无关，这正是它值钱的地方——方程再多，达到给定精度所需的投影次数不变，多出来的行只是让每次抽样有更多选择。

随机化治好了顺序依赖，没有治好病态。上面那张图的右半边，换成随机抽行仍然是右半边。

\subsection{它和随机梯度下降是同一件事的两种写法}

把每条方程写成一个损失项，

\[
f_i(x)=\frac{(a_i^{\trans}x-b_i)^2}{2\norm{a_i}_2^2},
\qquad
\nabla f_i(x)=\frac{a_i^{\trans}x-b_i}{\norm{a_i}_2^2}\,a_i ,
\]

那么 \(x-\eta\nabla f_i(x)\) 在 \(\eta=1\) 时正好是上面的投影公式。Kaczmarz 就是对 \(\frac1m\sum_i f_i\) 做步长为 1 的随机梯度下降，而按 \(\norm{a_i}_2^2/\norm{A}_F^2\) 抽样就是重要性采样。

这个对应不是类比，是等式，两边因此可以互相借经验。深度学习里每步只看一个 mini-batch 而不是全部样本，理由和这里一样：单步代价与数据总量无关，收敛按更新次数计而不是按遍历次数计。反过来，\hyperref[sec:09-Convex-Optimization/06-Optimization-Algorithms/01]{``梯度下降真正难在步长''}里最麻烦的那个问题——步长该取多大——在 Kaczmarz 这边消失了，因为二次损失沿单行方向的精确极小点有闭式解，\(\eta=1\) 就是精确线搜索的结果。这份便利有边界：只在每个 \(f_i\) 都是单行二次型时成立；一般损失函数没有这种结构，步长就得自己调。

从另一头看也有收获：随机梯度下降在病态问题上收敛慢，人们习惯归因于梯度噪声，而这张图给出了更朴素的解释——相邻两次更新的方向近乎相反，互相抵消。动量、预条件、行归一化，处理的都是同一个角落。

\subsection{行缩放改的是抽样，不是几何}

把某条方程两边同乘一个常数，超平面纹丝不动，投影结果也一模一样。但按行范数抽样时，它被选中的概率变了。实现里必须把``方程的几何''和``采样的策略''分开：前者对行缩放不变，后者不是。看到收敛行为随着数据的单位换算而改变，多半就是这两件事被混在了一起。

顺带一提，均匀抽样并不天然公平。范数很小的行同样携带一个完整的约束，只是它对 \(\norm{A}_F^2\) 的贡献小；若这条约束恰好是唯一能确定某个方向的，忽略它就等于放弃那个方向。另一个极端是每次都挑残差最大的行，这类似贪心坐标法，早期下降更快，但每步都要扫全部行，正好抵消了逐行方法的全部好处。比较时该看的是达到目标精度的总数据访问量，不是迭代计数。

\subsection{方程互相冲突时它不会自动给出最小二乘解}

有测量噪声时超平面没有公共交点，几张平面围出一个小区域。普通 Kaczmarz 会在这个区域里持续游走，落点取决于最后访问的是哪几行——它有极限集，但那个极限集一般不包含

\[
\argmin_x\norm{Ax-b}_2^2 .
\]

这是个容易踩的坑，因为迭代看上去很正常：残差降到某个水平后不再下降，曲线平了，像是收敛了。

想要最小二乘解就得改算法。随机 extended Kaczmarz 同时对列做投影，逐步剥掉 \(b\) 落在列空间之外的那部分，再对校正后的右端逐行投影；也可以把步长从 1 降下来并逐步衰减，退回成标准的 SGD 收敛模式；规模允许时，直接用 LSQR 这类基于 Krylov 的最小二乘算法更稳妥。选之前先问清楚目标：解一致方程、跟踪流式数据，还是最小化整体残差——这三件事需要三个不同的算法。

\subsection{整个单元其实只在回答一个问题}

回头看这一章的七篇，主线是清楚的：同一个 \(A^{-1}b\)，在计算机上要按矩阵的结构和数据的访问方式拆成完全不同的算法。有结构就用结构，一般方阵走带主元 LU，最小二乘走 QR，正定系统走 Cholesky；因子存不下就退到只需矩阵--向量乘的 Krylov，那时决定成败的从规模变成了谱；只要谱本身，就换成幂迭代与 QR 算法；只要主要方向，就用 SVD，同时接受秩在浮点下已经不是整数；连一次完整的矩阵--向量乘都做不起，就退到 Kaczmarz，用一行换一次投影。每退一步，需要的资源少一点，能拿到的保证也弱一点，而这个交换是自觉的，不是被迫的。

有一条线索贯穿全章：残差小从来不等于解对。它在正规方程里、在 Krylov 的停止准则里、在非正规矩阵的特征值里、在冲突方程的 Kaczmarz 里各出现了一次，形态不同，账是同一笔。

线性问题到此为止。下一章 \hyperref[ch:066]{``常微分方程数值方法：沿局部斜率走完整条轨迹''}换一类对象：未知的不再是一个向量，而是一整条随时间演化的轨迹，而唯一可用的信息是每一点上的变化率。稳定性这个词在那里会重新定义一次，但要问的问题还是老三样——花多少代价、准到什么程度、什么时候会崩。
